\section{问题分析}

\subsection{问题一的分析}

问题一要求在界面只有一次反射、透射的情形下建立确定外延层厚度的模型。入射光在外延层表面与衬底界面先后反射，两束光叠加形成干涉，反射率取决于两束光的相位差。难点在于相位差把折射率与厚度耦合在一起，两个量都未知。但条纹的幅度只由界面的反射与透射决定，条纹的位置只由光在外延层内的光程决定，因此可以先由峰谷反射率提取包络，按菲涅耳关系反推出折射率，再由相位确定厚度。折射率随波数变化，包络提取与相位反推都需逐峰谷对局部进行，两条路径分离又可相互检验。

\subsection{问题二的分析}

问题二要求把问题一的模型在实测光谱上实现，并分析结果的可靠性。数据的难点有三处：实测反射率含异常读数，峰谷识别受噪声干扰，碳化硅的折射率又随波数变化。先剔除异常点，再识别峰谷；每一对相邻峰谷按相位条件给出一个相互独立的厚度估计，多个估计取中位数作为结果，以离散程度衡量可靠性。折射率的色散使常数折射率近似只在部分谱段成立，需用文献名义值逐段核查，只让近似成立的谱段参与统计。可靠性还需要独立验证：同一晶圆片不同入射角的测量结果应当接近；再用全谱拟合从整条谱线的形状求厚度，与逐对峰谷的统计对照。几种结果彼此一致，厚度才可信。

\subsection{问题三的分析}

问题三的任务有两层：先推导多光束干涉产生的必要条件，再用它判别实测数据并计算厚度。光在膜内多次往返时，每一束后续光比前一束多经历一次界面反射与膜内衰减，衰减的快慢决定后续光束能否参与叠加；各束光要稳定干涉，相邻光束的光程差还须不超过光源的相干长度。判别的难点在于界面参数无法直接测量，且多光束干涉不改变条纹的周期、只改变条纹的形状，从周期上看不出来，只能从形状与幅度入手：振幅比可由峰谷反射率反推，条纹形状可由数据点相对中线的分布统计刻画，频域结构则反映谐波成分。一旦确认存在多光束干涉，包络线法失效，折射率需改取文献名义值，厚度可改由多光束模型全谱拟合、计入相移修正的峰谷法与傅里叶检验相互印证地求解；若判定不存在多光束干涉，问题二的算法仍然适用。
